\section{Absoluteness}

The question in this section is when a formula says the same thing in a class as it does in a bigger one.

\begin{boxdefinition}[Absoluteness]
    Let $M \ssq N$ be classes. We say a formula $\varphi\of{\bar{u}}$ is
    \begin{itemize}
        \item \textbf{downwards-absolute} from $N$ to $M$ iff for all $\bar{x} \in M^{<\omega}$ (matching $\bar{u}$ in length), $\varphi^N\of{\bar{x}} \to \varphi^M\of{\bar{x}}$.
        \item \textbf{upwards-absolute} from $M$ to $N$ iff for all $\bar{x} \in M^{<\omega}$, $\varphi^M\of{\bar{x}} \to \varphi^N\of{\bar{x}}$.
        \item \textbf{absolute} between $M$ and $N$ iff it is both, ie, iff for all $\bar{x} \in M^{<\omega}$, $\varphi^M\of{\bar{x}} \lr \varphi^N\of{\bar{x}}$.
    \end{itemize}
\end{boxdefinition}

\subsection{Absoluteness Theorems}

Every absoluteness theorem we prove will have an assumption that the smaller model (or sometimes, both models) are transitive. There are no absoluteness theorems in this course that do not mention transitivity, and it is a common mistake to try and apply them in cases where there isn't transitivity.

\begin{boxlemma}[$\Delta_0$-Absoluteness]
    Let $M$ and $N$ be classes\footnote{possibly proper classes, possibly sets}, with $M \ssq N$. Assume $M$ is transitive. Then, for every $\Delta_0$ formula $\varphi\of{\bar{u}}$ and $\bar{x} \in M^{<\omega}$ (matching $\bar{u}$ in length),
    \begin{align*}
        \varphi^M\of{\bar{x}} \lr \varphi^N\of{\bar{x}}
    \end{align*}
    That is, $\varphi\of{\bar{u}}$ is absolute between $M$ and $N$.
\end{boxlemma}
\begin{proof}
    We argue by induction on the complexity of $\varphi\of{\bar{u}}$.

    \begin{itemize}
        \item \underline{Atomic Case.} Note that
        \begin{align*}
            \parenth{u_0 = u_1}^M
            \qquad
            u_0 = u_1
            \qquad
            \parenth{u_0 = u_1}^N
        \end{align*}
        are all exactly the same. Similarly,
        \begin{align*}
            \parenth{u_0 \in u_1}^M
            \qquad
            u_0 \in u_1
            \qquad
            \parenth{u_0 \in u_1}^N
        \end{align*}
        are all exactly the same.

        \item \underline{Propositional Case.} If $\varphi$ is some propositional combination of $\varphi_0$ and $\varphi_1$, then the statement for $\vp$ is clear from the induction hypotheses for $\vp_0$ and $\vp_1$.

        \item \underline{Bounded Quantifiers.} Say $\vp\of{\bar{u}}$ is of the form $\exists v \in w \, \parenth{\psi\of{\bar{u}, v, w}}$. (This is enough, because bounded universal can be formulated as a negation of a bounded existential, and the propositional case covers negations.)

        Let $\bar{x}, z \in M$. Then, the following are equivalent:
        \begin{itemize}
            \item $\brac{\exists v \in z, \, \psi\of{\bar{x}, v, z}}^M$
            \item There is $y \in z \cap M$ such that $\psi^M\of{\bar{x}, y, z}$
            \item There is $y \in z$ such that $\psi^M\of{\bar{x}, y, z}$ (because $M$ is transitive, so $z = z \cap M$)
            \item There is $y \in z \cap N$ such that $\psi^M\of{\bar{x}, y, z}$ (because $z \ssq M \ssq N$, so $z = z \cap N$)
            \item There is $y \in z \cap N$ such that $\psi^N\of{\bar{x}, y, z}$
            \item $\brac{\exists v \in z, \, \psi\of{\bar{x}, v, z}}^N$
        \end{itemize}
    \end{itemize}
\end{proof}

Putting a block of unbounded quantifiers in front of a $\Delta_0$ formula costs us one of the two directions.

\begin{boxlemma}
    In the same situation as above, we have upward absoluteness from $M$ to $N$ for $\Sigma_1$ formulae and downward absoluteness from $N$ to $M$ for $\Pi_1$ formulae.
\end{boxlemma}

The same goes for formulae that are only $\Sigma_1$ or $\Pi_1$ up to provable equivalence in a theory, as long as both classes satisfy it.

% [CORRECTED] ``and non-empty'' is inserted --- the notes use ``theory'' for a set of axioms (``ZFC is an infinite theory''; ``For each axiom $\sigma \in \ZF$ (\textbf{including} foundation), $\sigma^{\WF}$''), so for a class, ``model of $T$'' reads as $\sigma^M$ for every axiom $\sigma$ of $T$, which $M = \emptyset$ satisfies whenever every axiom of $T$ begins with $\forall$. Then $\exists u \parenth{u = u}$ is $\Delta_1^T$, being logically equivalent to $\forall u \parenth{u = u}$, but it is true in any non-empty $N$ and false in $M$: the $T$-equivalence only carries over to a non-empty class, as in the proof-transfer lemma of the section on relative consistency. If ``model of $T$'' is instead read as $\theta^M$ for every \textit{theorem} $\theta$ of $T$, or, for a set $M$, as $\parenth{M, \in} \models T$, whose universe is non-empty by definition, then $M \neq \emptyset$ already and the insertion is redundant, though harmless
\begin{boxlemma}
    Let $T$ be a theory and $M \ssq N$ be models of $T$ (with $M$ transitive and non-empty). Then,
    \begin{itemize}
        \item $\Sigma_1^T$ formulae are upwards-absolute
        \item $\Pi_1^T$ formulae are downwards-absolute
        \item $\Delta_1^T$ formulae are absolute
    \end{itemize}
\end{boxlemma}

\subsection{Absoluteness of ZFC Axioms}

For this discussion, fix a class $M$. Let's list out the axioms of $\ZFwoF$, relativized to $M$.

We begin with Extensionality, which relativized to $M$ reads
\begin{align*}
    \brac{\forall A \forall B \brac{\forall x \parenth{x \in A \lr x \in B} \to A = B}}^M
\end{align*}
Or, equivalently,
\begin{align*}
    \brac{\forall A \forall B \brac{\parenth{A \ssq B \land B \ssq A} \to A = B}}^M
\end{align*}
% [CORRECTED] was: ``Indeed, all sub-formulae of the above are $\Delta_0$.'' --- both displays begin with the unbounded $\forall A \forall B$, and the first also has an unbounded $\forall x$, so read as a claim about them it is false; what is $\Delta_0$, and what the lemma below uses, is the formula under the two leading quantifiers
where $A \ssq B$ is the $\Delta_0$ formula $\parenth{\forall x \in A}\parenth{x \in B}$. Indeed, all sub-formulae of $\parenth{A \ssq B \land B \ssq A} \to A = B$ are $\Delta_0$.

\begin{boxlemma}
    If $M$ is transitive, then Extensionality$^M$.
\end{boxlemma}
\begin{proof}
    By $\Delta_0$-absoluteness, for $A, B \in M$ the formula $\parenth{A \ssq B \land B \ssq A} \to A = B$ holds in $M$ iff it holds in $V$. Extensionality says it holds in $V$ for all $A$ and $B$, so in particular it holds in $M$ for all $A, B \in M$, which is what Extensionality$^M$ says.
\end{proof}

Assume (here on after) that \underline{$M$ is transitive}.

% [FILLED] this lemma and its proof are supplied; the lecture's notes stop at the transitivity assumption
Next comes Comprehension. For a formula $\varphi\of{x, \bar{y}}$, the corresponding instance of Comprehension, relativized to $M$, reads
\begin{align*}
    \forall A \in M \, \forall \bar{y} \in M \, \exists B \in M \, \forall x \in M \brac{x \in B \lr \parenth{x \in A \land \varphi^M\of{x, \bar{y}}}}
\end{align*}
and transitivity lets us say exactly when this holds.

\begin{boxlemma}[Scheme] \label{Ch3:Lemma:Comprehension_In_M}
    For each formula $\varphi\of{x, \bar{y}}$, the instance of Comprehension for $\varphi$ holds in $M$ iff
    \begin{align*}
        \setst{x \in A}{\varphi^M\of{x, \bar{y}}} \in M
    \end{align*}
    for all $A, \bar{y} \in M$. In particular, Comprehension$^M$ holds if every subset of a member of $M$ is itself a member of $M$.
\end{boxlemma}
\begin{proof}
    Fix $A, \bar{y} \in M$. Since $M$ is transitive, $A$ and every $B \in M$ are subsets of $M$, so for $x \notin M$ both sides of the biconditional above are false, and the innermost $\forall x \in M$ may as well range over every $x$. So what is being asked of $B$ is exactly that $B = \setst{x \in A}{\varphi^M\of{x, \bar{y}}}$, and there is such a $B$ in $M$ iff this set is a member of $M$. It is a subset of $A$, which gives the last part.
\end{proof}

From now on, assume Comprehension$^M$.

% [FILLED] this lemma and its proof are supplied; the lecture's notes have only the heading
Pairing is simpler, because atomic formulae relativize to themselves: Pairing$^M$ reads
\begin{align*}
    \forall x \in M \, \forall y \in M \, \exists A \in M \, \parenth{x \in A \land y \in A}
\end{align*}
ie, any two members of $M$ lie in a common member of $M$. With Comprehension$^M$ in hand, we can cut that common member down to size.

\begin{boxlemma}
    Pairing$^M$ holds iff $\set{x, y} \in M$ for all $x, y \in M$.
\end{boxlemma}
\begin{proof}
    If $\set{x, y} \in M$ whenever $x, y \in M$, then Pairing$^M$ certainly holds. Conversely, given $x, y \in A \in M$, \Cref{Ch3:Lemma:Comprehension_In_M} gives $\setst{z \in A}{z = x \lor z = y} \in M$, since $z = x \lor z = y$ is its own relativization, and that set is $\set{x, y}$.
\end{proof}
